% Part: turing-machines
% Chapter: undecidability
% Section: unsolvability-decision-problem

\documentclass[../../../include/open-logic-section]{subfiles}

\begin{document}

\olfileid{tur}{und}{dec}
\olsection{Masalah Putusan}

Kita nyebut logika tataran kapisan \emph{bisa diputusake}
yen lan mung yen ana cara efektif kanggo nemtokake
apa sawijining !!{sentence} kang diwenehake sah
utawa ora. Pranyata cara kaya mangkono ora ana:
masalah nemtokake kasahihan ukara tataran kapisan
ora bisa dirampungake.

Kanggo netepake asil negatif kang wigati iki,
kita mbuktekake yen masalah putusan ora bisa
dirampungake dening mesin Turing. Tegese,
ora ana mesin Turing kang, saben diwiwiti
ing pita kang ngemot !!{sentence} tataran kapisan,
mesthi mandheg lan ngasilake $1$ utawa~$0$
miturut apa !!{sentence} iku sah utawa ora.
Miturut tézis Church--Turing, saben fungsi
kang bisa diitung uga bisa diitung dening
mesin Turing. Mula yen ``fungsi kasahihan''
iki bisa diitung kanthi efektif, fungsi iku
uga bisa diitung dening mesin Turing.
Yen ora bisa diitung dening mesin Turing,
fungsi iku uga ora bisa diitung kanthi efektif.

Strategi kita kanggo mbuktekake yen masalah
putusan ora bisa dirampungake yaiku ngreduksi
masalah mandheg menyang masalah iku. Tegese:
kita wis mbuktekake yen fungsi~$h(e,w)$,
kang menehi output~$1$ yen mesin Turing
kang diterangake dening~$e$ mandheg kanggo
input~$w$ lan output~$0$ yen ora, ora
bisa diitung dening mesin Turing. Kita
bakal nuduhake yen ana mesin Turing kang
mutusake kasahihan ukara tataran kapisan,
mula uga ana mesin Turing kang ngitung~$h$.
Amarga $h$ ora bisa diitung dening mesin
Turing, mesin Turing kang mutusake kasahihan
uga ora mungkin ana.

Langkah kapisan ing strategi iki yaiku nuduhake
yen kanggo saben input~$w$ lan mesin Turing~$M$,
kita bisa kanthi efektif njlentrehake
!!a{sentence} $!T(M, w)$ kang nggambarake
himpunan instruksi~$M$ lan input~$w$,
lan !!a{sentence}~$!E(M, w)$ kang nyatakake
``$M$ pungkasane mandheg'' saengga:
\begin{quote}
  $\Entails !T(M, w) \lif !E(M,w)$ yen lan mung yen $M$ mandheg kanggo input~$w$.
\end{quote}
Bagean utama bukti kita yaiku njlentrehake
ukara loro iki, $!T(M, w)$ lan~$!E(M, w)$,
banjur mbuktekake yen $!T(M, w) \lif !E(M, w)$
sah yen lan mung yen $M$~mandheg kanggo input~$w$.

\end{document}
